\section{Quality Changes Are Not Gap Recovery}
\label{app:quality}

\paragraph{Matched protocol Q3 (V46/V54).}
Qwen2.5-7B-Instruct and Llama-3.1-8B-Instruct each have three independently
trained top-eight-layer checkpoints, paired with their base models under full,
DDI-QF and blind masks. Each seed/dataset has 100 common natural-QA IDs across
architectures and all six arms; the synthetic control has 20 records. Greedy
generation is capped at 64 tokens, batch one, bf16. Retained prediction
provenance records A100 80GB PCIe, torch 2.13.0+cu126, code/config identity,
seed and checkpoint directory. These are three selected instances, not a
population sample of training outcomes. Per-row provenance, rather than a
single campaign-wide clean-source assertion, is the authority.

Let $B_F,B_D,B_B$ denote base full/DDI/blind means and $T_F,T_D,T_B$ tuned
means. The preregistered statistic
\begin{equation}
C=1-\frac{T_F-T_D}{B_F-B_B}
\end{equation}
normalizes the remaining tuned gap by base document headroom. It is not the
fraction of the original full--DDI gap recovered through training. A decrease
in full-path quality can raise $C$ without any DDI gain. V54 therefore retains
$C$ and its original gates but separately reports
\begin{equation}
\Delta_D=T_D-B_D,\quad\Delta_F=T_F-B_F,\quad
\Delta_G=(T_F-T_D)-(B_F-B_D)=\Delta_F-\Delta_D.
\end{equation}
Positive $\Delta_D,\Delta_F$ mean improvement; negative $\Delta_G$ means gap
narrowing. Loss is sign-oriented before contrasts; stored raw loss remains
lower-is-better. Presence and loss are supporting controls, not substitutes for
failed primary EM/F1 gates.

\begin{figure}[H]
\centering% Generated by paper/render.py; aggregates: paper/evidence.json.
% Sources: experiments/transfer_v54_posthoc_sensitivity.json sha256=7da99a6fcb70f5c3bbf88fd13ca71e3c4e2182c6daf381f97a90d9fe3ec5f863
\begin{tikzpicture}
\begin{groupplot}[group style={group size=2 by 1,horizontal sep=1.55cm},width=.455\linewidth,height=5.2cm,grid=major,grid style={gray!15},tick label style={font=\small},label style={font=\small},title style={font=\small\bfseries},legend style={font=\small,draw=none,at={(.5,-.28)},anchor=north,legend columns=2,/tikz/every even column/.append style={column sep=5pt}}]
\nextgroupplot[title={Qwen2.5-7B},ylabel={paired change},xtick={0,1,2,3},xticklabels={H-EM,H-F1,M-EM,M-F1},ymin=-.18,ymax=.17,ybar,bar width=5pt,enlarge x limits=.18,legend columns=3]
\addplot[fill=ddiBlue,draw=none] coordinates {(0,0.130000) (1,0.134740) (2,0.066667) (3,0.080792)};
\addplot[fill=ddiGold,draw=none] coordinates {(0,-0.026667) (1,-0.018846) (2,-0.033333) (3,-0.002098)};
\addplot[fill=ddiPurple,draw=none] coordinates {(0,-0.156667) (1,-0.153586) (2,-0.100000) (3,-0.082890)};
\legend{DDI,Full,Gap}
\nextgroupplot[title={Llama-3.1-8B},ylabel={paired change},xtick={0,1,2,3},xticklabels={H-EM,H-F1,M-EM,M-F1},ymin=-.18,ymax=.17,ybar,bar width=5pt,enlarge x limits=.18,legend columns=3]
\addplot[fill=ddiBlue,draw=none] coordinates {(0,0.000000) (1,0.000012) (2,0.033333) (3,0.038378)};
\addplot[fill=ddiGold,draw=none] coordinates {(0,-0.053333) (1,-0.025424) (2,0.020000) (3,0.032865)};
\addplot[fill=ddiPurple,draw=none] coordinates {(0,-0.053333) (1,-0.025436) (2,-0.013333) (3,-0.005513)};
\legend{DDI,Full,Gap}
\end{groupplot}
\end{tikzpicture}

\caption{Equal-seed means of paired training changes. H/M denote Hotpot/MuSiQue.
Qwen's positive DDI changes coexist with negative full-path changes. Llama's
Hotpot gap narrowing largely reflects full-path degradation. Bars are point
estimates, not an architecture-equivalence or parity test.}
\end{figure}

Neither architecture reaches natural-QA parity under the retained rules. The
only valid original architecture contrast is MuSiQue F1: Llama minus Qwen
$C=+0.028$, equal-seed interval $[-0.085,+0.179]$. This does not establish a
practical architecture difference or equivalence. Blind-control movement
invalidates Llama Hotpot; insufficient headroom/preservation invalidates other
specified cells. Small remaining gaps alone cannot override those gates.

V54's post-hoc sensitivity jointly resamples common record IDs across the six
arms and both architectures within each seed (10,000 draws, RNG seed 54).
It does not resample a population of training seeds. Conditional intervals for
Qwen DDI F1 change are $[+.095,+.176]$ on Hotpot and $[+.042,+.120]$ on MuSiQue.
They clarify the observed training changes but do not replace the original
decision procedure. Full results and inputs are in
\path{experiments/transfer_v54_posthoc_sensitivity.json} and
\path{experiments/transfer_v54_inputs.json}; see \path{VIBEQUEUE-54-REPORT.md}.

\newpage
\section{Cache Anatomy and Provenance}
\label{app:cache}

\paragraph{Request-level derivation.}
The postprocessor reads the hash-pinned V53 compact summary, 540 timing rows
and 192 validation rows. It requires all $(\mathrm{seed},\mathrm{request},
\mathrm{mode})$ keys, exact five-by-500 document geometry, query length 32,
16 generated IDs and 15 decode forwards. Measured requests are indices 10--59;
the corresponding validation case is \texttt{stream\_k}, without renumbering.
Cached/fresh generated IDs must match those retained in validation. The
renderer recomputes each QPS as $50/\sum t_i$ and checks the compact summary.
Grouping uses hits in the \emph{timing} cache, not the separate validation cache.

\begin{table}[H]
\centering\small% Generated by paper/render.py; aggregates: paper/evidence.json.
% Sources: VIBEQUEUE-53-REPORT.md sha256=7cd9223e128c5d82b61e59575035fe4bb54c9f48c5da721723fd5d214d85195e; results/cache_v53_local_20260907/summary.json sha256=28ae886d9dae851ed090b7ebb0744b33926060d9499c3f16a01221810a0e2189; results/cache_v53_local_20260907/timing.jsonl sha256=055cf47312269ca147c6d4517d26b82aabc910b11fd6ca07f5eba5e3f54153a8; results/cache_v53_local_20260907/validation.jsonl sha256=8597c0b5394b4b17248065ad1ebc386a0a2d864f089f8e65d51d5e293f7b8d4a
\resizebox{\linewidth}{!}{%
\begin{tabular}{rrrrrrr}
\toprule
Hits & $n$ & Seed counts (0/1/2) & Equal IDs & Cache ms & Fresh ms & Paired ratio \\
\midrule
3 & 4 & 1/1/2 & 4/4 & 729.8 & 627.1 & 0.865 \\
4 & 33 & 11/12/10 & 28/33 & 606.5 & 625.4 & 1.034 \\
5 & 113 & 38/37/38 & 99/113 & 498.3 & 626.8 & 1.259 \\
\bottomrule
\end{tabular}}

\caption{Post-hoc measured hit bins. Cache/fresh milliseconds are medians on
matched requests; the last column is the median of within-request fresh/cache
latency ratios, not the ratio of medians. Equal IDs counts exact 16-token
sequence agreement. Small bins and unequal seed composition preclude a causal
hit-rate effect. No request-level confidence intervals are asserted.}
\end{table}

\paragraph{Behavioral anatomy.}
Each seed validates cold, warm, permuted and partial-reuse cases plus all 60
stream requests. Greedy failures are 0, 6 and 19 of 64 comparisons by seed;
measured-stream failures are 0, 4 and 15 of 50. The additional failures comprise
five warmup requests and one explicit permutation. Every numerical comparison
fails the unchanged $10^{-4}$ absolute/relative gate. All 384 natural-generation
outputs reach the 16-token cap; no early EOS was observed in this exact-ID stream.
Natural checks reuse already populated validation entries and are not cold
latency measurements. Teacher-forced comparisons follow fresh DDI tokens to
keep same-prefix logit errors meaningful after greedy trajectories diverge.

For example, seed 1 \texttt{stream\_7} differs at its first token: fresh chooses
220, cache 198. The fresh top-two margin is .125, the cache margin zero, and
maximum first-step logit error .21875. This is consistent with argmax changes
near a tie; it does not identify a kernel-level root cause or prove an adequate
precision fix. A passing tiny fp32 CPU test cannot establish full-model bf16
equivalence. These examples are diagnostic, not a QA correctness judgment.

\paragraph{Three generations must remain separate.}
The pre-V51 1.72$\times$ claim is withdrawn because the benchmark populated
but did not consume real KVs. V51's real-cache 1.219$\times$ claim is also
withdrawn: random IDs were decoded and retokenized, changing nominal lengths,
and full versus DDI decode work was not contractually equal. V53 is a new
exact-ID diagnostic on a different torch runtime, not an identical replication.
Its seed-wise rates must not be pooled with either older campaign.

\paragraph{Source identity and scope.}
The V53 manifest hashes checkpoint/tokenizer contents and preserves executing
sources. Base SHA is \texttt{16dd67d7e53a8168cf6580e6d0197c219e6ba198}; the
diagnostic script was untracked at launch, despite an empty tracked diff.
The manifest says \texttt{clean\_revision=false}, and both manifest and summary
say \texttt{admissible\_paper\_evidence=false}. The report remains unchanged:
\path{VIBEQUEUE-53-REPORT.md}. Exact inputs, logits, source snapshots and audit
live under \path{results/cache_v53_local_20260907/}. This paper additionally pins
the compact summary/timing/validation byte sizes and full SHA-256 values in
\path{paper/evidence.json}; hashes make postprocessing reproducible, not the
execution clean or correctness-preserving.

\paragraph{Where the optimized warm-path win came from (V69).}
The audited all-hit reuse acceleration of the main text (221.24\,ms before the
V65 optimizations, 62.30\,ms after, 3.55$\times$) is attributed by a
derivation over pinned artifacts, not by a new run: the pre/post totals trace
to the confirmation-frozen V60/V65 runs (V64-adjudicated) and the phase leaves
are uniformly rescaled from the V65 overhead profile (factor 0.945248).
Store-time validation markers remove 85.61\,ms (53.9\% of the reduction) and
the 2\,GiB LRU resident tier removes the synchronous per-document disk load of
59.14\,ms (37.2\%); the remaining 14.20\,ms (8.9\%) is a bounded envelope
within the pre-optimization cache\_assemble budget --- inferred, not directly
measured --- and the unattributed wall remainder is 3.48\,ms (1.6\% of the
pre-total, within the declared 5\% tolerance). Development-role evidence; it
supports no confirmation-quality claim.

\begin{table}[H]
\centering\small% Generated by paper/render.py from paper/evidence.json.
% Sources: results/v69_attribution_20260910T162613Z/attribution.json sha256=90f34eec2d9b24426a0af79ba757bedad178b8cb1abecf503cb0a87aea2188d1.
\begin{tabular}{llrr}
\toprule
Optimization (V65) & Pre-optimization phase removed/reduced & ms removed & \% of reduction \\
\midrule
OPT-1 store-time validation markers & validate\_entry (per-request geometry+isfinite scan) & 85.61 & 53.9 \\
OPT-2 2 GiB LRU resident tier & disk\_load (synchronous per-doc torch.load) & 59.14 & 37.2 \\
OPT-3 single staged pinned H2D copy per layer (bounded envelope) & cache\_assemble (per-entry .to(device) + cat) & 14.20 & 8.9 \\
\midrule
Total removed & & 158.94 & 100.0 \\
\midrule
\multicolumn{4}{l}{Pre-optimization all-hit reuse 221.24 ms $\to$ post-optimization 62.30 ms (3.55$\times$).} \\
\multicolumn{4}{l}{OPT-3 is a bounded envelope, not a direct measurement: it fits within the 23.04 ms pre-optimization} \\
\multicolumn{4}{l}{cache\_assemble budget; the implied post-optimization assemble of 8.84 ms is inferred, not measured.} \\
\multicolumn{4}{l}{Unattributed wall remainder 3.48 ms (1.6\% of the pre-total, within the declared 5\% tolerance).} \\
\bottomrule
\end{tabular}

\caption{Attribution of the 221.2\,ms $\to$ 62.3\,ms all-hit reuse reduction
(V69 derivation over pinned V60/V65 artifacts). OPT-1 and OPT-2 are measured
pre-optimization phases removed by the V65 changes; OPT-3 is a bounded
envelope, not a direct measurement. Development-role evidence; no
confirmation-quality claim.}
\label{tab:v69attribution}
\end{table}

\newpage
\section{Artifact Ledger and Missing Joint Design}
\label{app:ledger}

Paths below are relative to the repository root. The authoritative full hashes
and byte sizes are in \path{paper/evidence.json}. No report was modified to
retroactively admit a result. \path{paper/generated/analysis.json} contains the
new request-bin statistics, seed QPS and retained failure totals.

\begin{small}
\begin{tabular}{@{}p{2.2cm}p{8cm}p{5.2cm}@{}}
\toprule
Evidence & Exact artifacts & Row selection and unit \\
\midrule
Q1 / T2 & \path{results/runs.jsonl}; \path{results/sprint_a_strict_base.jsonl} &
\texttt{phase2\_train\_free}, Qwen-7B; tags \texttt{88162acd}/\texttt{fee7d916}; 100 IDs per arm. \\[6pt]
Q2 / T5 & \path{results/phase3_n2.jsonl}, \path{results/phase3_n5.jsonl}, \path{results/phase3_n10.jsonl} &
Hotpot full/QF, seed 0, tag \texttt{3a59828b}; match final ID suffix. \\[6pt]
P1 / T3 & \path{results/runs.jsonl} &
\texttt{phase3\_timing}, Qwen-7B, \texttt{L=512:Q=64}, SDPA, tag \texttt{88162acd}. \\[6pt]
Q3 / V54 & \path{experiments/transfer_v54_inputs.json}; \path{experiments/transfer_v54_posthoc_sensitivity.json} &
12 prediction paths; pair dataset/seed/record ID across six arms and architectures. \\[6pt]
C1 / V53 & \path{results/cache_v53_local_20260907/summary.json}, \path{timing.jsonl}, \path{validation.jsonl} in that directory &
Timing: \texttt{warmup=false}; join seed/request to validation \texttt{stream\_k}. \\
\bottomrule
\end{tabular}
\end{small}

\paragraph{Checkpoint identity cannot be inferred from the display label.}
Q2's raw run IDs name the base model, but
\path{scripts/phase3_doc_sweep2.sh} and
\path{logs/phase3_doc_sweep2.log} identify the checkpoint
\path{results/checkpoints/Qwen__Qwen2.5-7B-Instruct__ddi_qf__8L}.
File-specific Q2 config hashes are \texttt{94b5cfd00bbc44e0},
\texttt{59a0f85d992e8ca4}, and \texttt{b5f6c833f633be8e}. Do not aggregate all
base-labeled rows in the shared historical store. Q3 instead records
\texttt{checkpoint\_revision} and separate base/tuned identities in predictions.
Its exact per-seed manifests are under \path{results/transfer_heldout/}.

\paragraph{Schemas and missing fields.}
Legacy scalar JSONL has \texttt{run\_id, metric, value, provenance}; quality
record IDs are the run-ID suffix. Modern predictions have a separate
\texttt{record\_id}, decoded prediction, possible answers, a four-metric object,
and provenance. Neither QA schema stores per-record prefill times, document
token lengths or cache hit states. Timing repeats are JSON strings inside the
scalar \texttt{value}. V53 has request-level geometry, IDs, hits and latency,
but no gold QA answers. These schema differences explain the missing cells in
Figure~\ref{fig:coverage}; matching only the model and mask is not enough.

\paragraph{Historical Flex lineage.}
Figure~\ref{fig:flexhistory} parses detailed milliseconds, not rounded ratios,
from hash-pinned \path{logs/timing_flex2.log} and
\path{logs/timing_flex_model3.log}. The scripts preserved at Git revision
\texttt{2029129}, \path{scripts/timing_flex_test.py} and
\path{scripts/timing_flex_model.py}, distinguish raw Q/K/V timing from the
whole-model forward. Their document mask lacks an intra-document causal
condition. Preserved logs establish reported timings, not exact launch-source
identity or current-mask parity. The separate 3.98$\times$ archived claim is
recorded in that revision's changelog; absent its matching raw log, it remains
unplotted. The old square-matrix sparsity summaries are not reused.

\paragraph{Minimum future joint record.}
Retain an immutable request ID, exact document/query IDs, gold answers,
checkpoint/tokenizer/source hashes, mask/layout, cache namespace and pre-state,
hit identities, actual geometry, generated IDs and stopping reason. Record
natural-generation EM/F1 separately from fixed-work TTFT, end-to-end time and
decode forwards. Randomize or counterbalance execution order and preserve
per-seed pairing. Evaluate cold/warm/reordered/partial states and a real
optimized prefix-cache baseline before claiming an advantage over prefix reuse.

Reports \path{VIBEQUEUE-54-REPORT.md}, \path{VIBEQUEUE-53-REPORT.md},
\path{VIBEQUEUE-55-REPORT.md} and \path{EVIDENCE_TO_CLAIM.md} define the retained
estimands, withdrawals and failed kernel gates. A future matched campaign may
fill the coverage map; the present appendix does not count that unrun design
as an experiment or relax any original gate.
